\centerline{\bf \Huge Солянка I}
\centerline{\bf \Huge }

\begin{flushright}
        \scriptsize{}
\end{flushright}

\problem{1}
По кругу стоят 200 героев из 10 различных фракций, в каждой из которых 20 героев. Докажите, что в каждой группе можно выбрать вожака таким образом, чтобы никакие два вожака не стояли бы рядом.

\problem{2}
Отображение $f(x)$ называется \textit{полиномиальным}, если оно задаётся каким-то многочленом с действительными коэффициентами. Число $x$ называется \textit{неподвижной точкой} этого отображения, если $f(x) = x$. Оказалось, что у отображения полиномиального отображения $f(x)$ нет неподвижных точек. Докажите, что у $f(f(x))$ так же нет неподвижных точек.

\problem{3}
В треугольнике $A B C$ точка $O$ --- центр описанной окружности, $A D$ --- высота, опущенная на сторону $B C$. Пусть $E$ и $F$ --- проекции точек $B$ и $C$ на прямую $A O$. Обозначим через $G$ точку пересечения прямых $A C$ и $D E$, а через $H$ --- точку пересечения прямых $A B$ и $D F$. Докажите, что точки $A, G, D, H$ лежат на одной окружности.

\problem{4}
Существует ли такое натуральное $n$, что число $3^n+2 \cdot 17^n$ является точным квадратом?

\problem{5}
Существует ли такой квадратный трехчлен $f(x)$, что $[f(x)]=f([x])$ ? (Здесь $[a]$ целая часть числа $a$, т.е. наибольшее целое число, не превосходящее $a$.)

\problem{6}
Вписанная окружность $\gamma$ треугольника $A B C$ касается его стороны $B C$ в точке $D$. Обозначим центр вневписанной в угол $A$ окружности через $I_a$, а середину отрезка $D I_a$ через $M$. Докажите, что отрезок касательной, проведённой из точки $M$ к $\gamma$, равен $M B$.